% Supplementary-information stub text, generated 2026-09-04 from the pipeline
% (see journal/notes/paper_technical_check_20260904.md for the source of every
% number). Every paragraph is tagged [AI generated]; delete the tag once a human
% has rewritten and checked it. Fragment to \input into the Overleaf SI; Table S1
% needs \usepackage{booktabs} and the file paper/tables/out/datasets_table.tex.

\section*{S1. Datasets}

[AI generated] Table~\ref{tab:si-datasets} lists every dataset in its analysis
sample, that is, after all exclusions described below: the five Wordbank
longitudinal datasets that enter the imputed-input fits and the longitudinal
demographic arm; the six English datasets with input or processing measures; and
the 31 Wordbank cross-sectional languages fit for the demographic analysis, of
which 28 (sex) and 12 (maternal education) pass the display threshold.

\begin{table}[htbp]
\centering\footnotesize
\caption{[AI generated] Datasets and analysis samples. Children and administrations
are counts after exclusions; ages are administration ages in months; recordings are
day-long LENA recordings or BabyView videos entering the input measurement model;
LWL sessions are looking-while-listening administrations entering the processing
measurement model. Regenerated from the fitted bundles by
\texttt{paper/tables/datasets\_table.R}.}
\label{tab:si-datasets}
% Generated by paper/tables/datasets_table.R from paper/cache/table1_datasets.csv -- do not edit by hand.
\begin{tabular}{llrrrrlrr}
\toprule
Language & Source & $N$ & Admin. & Ages (mo) & Mean & Input & Rec. & LWL \\
\midrule
\addlinespace[2pt]\multicolumn{9}{l}{\textit{Wordbank longitudinal (imputed-input fits; demographic BLUP arm)}} \\[1pt]
English (American) & Thal et al.\ [year TODO] & 653 & 1,947 & 12--29 & 19.1 & -- & -- & -- \\
English (American) & Smith et al.\ [year TODO] & 316 & 935 & 16--30 & 22.1 & -- & -- & -- \\
English (American) & Marchman et al.\ [year TODO] & 2,136 & 4,447 & 8--30 & 19.0 & -- & -- & -- \\
Norwegian & Simonsen et al. (2014) & 1,630 & 6,446 & 8--36 & 24.9 & -- & -- & -- \\
Japanese & Hagihara et al. (2023) & 187 & 471 & 8--34 & 17.4 & -- & -- & -- \\
\addlinespace[2pt]\multicolumn{9}{l}{\textit{Input / processing (joint and separate measurement-model fits)}} \\[1pt]
English (American) & Long et al. (2024) & 22 & 116 & 9--32 & 19.1 & head-cam & 5,739 & -- \\
English (American) & Egan-Dailey \& Bergelson (2025) & 44 & 514 & 6--18 & 12.2 & LENA & 525 & 256 \\
English (American) & Adams et al. (2018) & 67 & 167 & 13--27 & 20.1 & LENA & 126 & 440 \\
English (American) & Fernald et al. (2013) & 105 & 276 & 16--30 & 21.7 & LENA & 157 & 134 \\
English (American) & Fernald \& Marchman (2012) & 119 & 382 & 14--30 & 22.9 & -- & -- & 419 \\
English (American) & Fernald, Perfors \& Marchman (2006) & 56 & 160 & 18--25 & 21.3 & -- & -- & 125 \\
\addlinespace[2pt]\multicolumn{9}{l}{\textit{Wordbank cross-sectional (Frank et al., 2017; one administration per child)}} \\[1pt]
Arabic (Saudi) & Wordbank & 258 & 258 & 8--16 & 11.9 & -- & -- & -- \\
Cantonese & Wordbank & 1,208 & 1,208 & 16--30 & 23.4 & -- & -- & -- \\
Catalan & Wordbank & 1,436 & 1,436 & 8--30 & 19.7 & -- & -- & -- \\
Croatian & Wordbank & 627 & 627 & 8--30 & 19.0 & -- & -- & -- \\
Czech & Wordbank & 493 & 493 & 16--30 & 23.0 & -- & -- & -- \\
Danish & Wordbank & 6,112 & 6,112 & 8--36 & 20.8 & -- & -- & -- \\
Dutch & Wordbank & 446 & 446 & 8--31 & 15.6 & -- & -- & -- \\
English (American) & Wordbank & 8,685 & 8,685 & 8--30 & 19.2 & -- & -- & -- \\
English (Australian) & Wordbank & 1,493 & 1,493 & 12--30 & 20.6 & -- & -- & -- \\
Estonian & Wordbank & 1,235 & 1,235 & 16--30 & 23.3 & -- & -- & -- \\
Finnish & Wordbank & 234 & 234 & 9--24 & 16.5 & -- & -- & -- \\
French (French) & Wordbank & 1,290 & 1,290 & 8--30 & 21.6 & -- & -- & -- \\
German & Wordbank & 1,181 & 1,181 & 18--30 & 24.1 & -- & -- & -- \\
Hebrew & Wordbank & 534 & 534 & 11--35 & 27.4 & -- & -- & -- \\
Hungarian & Wordbank & 369 & 369 & 16--30 & 23.1 & -- & -- & -- \\
Italian & Wordbank & 1,400 & 1,400 & 7--36 & 20.4 & -- & -- & -- \\
Japanese & Wordbank & 846 & 846 & 8--36 & 20.4 & -- & -- & -- \\
Kigiriama & Wordbank & 205 & 205 & 8--30 & 16.1 & -- & -- & -- \\
Korean & Wordbank & 2,229 & 2,229 & 8--36 & 22.9 & -- & -- & -- \\
Latvian & Wordbank & 683 & 683 & 8--36 & 22.3 & -- & -- & -- \\
Mandarin (Beijing) & Wordbank & 1,056 & 1,056 & 16--30 & 23.0 & -- & -- & -- \\
Mandarin (Taiwanese) & Wordbank & 2,654 & 2,654 & 8--36 & 22.0 & -- & -- & -- \\
Norwegian & Wordbank & 7,358 & 7,358 & 8--36 & 23.3 & -- & -- & -- \\
Portuguese (European) & Wordbank & 4,322 & 4,322 & 8--30 & 19.7 & -- & -- & -- \\
Russian & Wordbank & 1,803 & 1,803 & 8--36 & 21.4 & -- & -- & -- \\
Slovak & Wordbank & 1,712 & 1,712 & 8--36 & 20.0 & -- & -- & -- \\
Spanish (Argentinian) & Wordbank & 783 & 783 & 16--33 & 23.3 & -- & -- & -- \\
Spanish (European) & Wordbank & 1,005 & 1,005 & 8--30 & 18.4 & -- & -- & -- \\
Spanish (Mexican) & Wordbank & 1,915 & 1,915 & 8--30 & 19.1 & -- & -- & -- \\
Swedish & Wordbank & 1,358 & 1,358 & 8--28 & 19.0 & -- & -- & -- \\
Turkish & Wordbank & 3,537 & 3,537 & 8--36 & 21.4 & -- & -- & -- \\
\bottomrule
\end{tabular}

\end{table}

\subsection*{S1.1 Wordbank data}

[AI generated] MB-CDI data were downloaded with the \texttt{wordbankr} package
(version 2.0.0) from the Wordbank v2.0 release hosted on Redivis. We used Words \&
Gestures (WG) and Words \& Sentences (WS) forms, including their short forms, and
coded production from the instrument's logical \texttt{produces} field.

[AI generated] For the longitudinal datasets, a child was identified by the
combination of dataset and study-internal participant identifier, which links a
child's WG and WS administrations; Wordbank's own child identifier fails to do so
in some datasets and would split cross-form children into two. Items were
harmonized across forms by their cross-linguistic \texttt{uni\_lemma} wherever the
mapping is unambiguous within a form, so that a child's WG and WS administrations
are scored on one item set. A child entered the longitudinal sample with at least
two administrations at distinct ages.

[AI generated] Trajectory quality control removed administrations that produce
implausible ``craters'' in a child's proportion of words produced: an
administration was dropped when a later administration fell below the child's
earlier peak by more than 25\% of that peak and by at least 0.05, with the peak
itself at least 0.10, together with an analogous rule for implausible jumps
[check exact jump rule in \texttt{qc\_clean\_child()}]. This removed 90 Marchman
and 50 Norwegian administrations.

[AI generated] We restricted Wordbank to monolingual, typically developing
samples by excluding datasets whose origin is labelled bilingual and the two
clinical datasets (Edgin; Byers). The English longitudinal sample pools the Thal
(653 children), Smith (316) and Marchman (2,136) datasets into one dataset of
3,105 children, 7,329 administrations and 681 harmonized items, with the anchor
age $a_0$ set to the median administration age (18 months); Norwegian comprises
1,630 children and 6,446 administrations. Japanese (187 children) enters only the
longitudinal demographic arm. Smith carries no demographic information, and Thal
and Japanese carry sex but not caregiver education.

\subsection*{S1.2 Input and processing datasets}

[AI generated] Vocabulary data for the six English datasets come from item-level
MB-CDI files: the Stanford Language Learning Lab files for Adams et al.\ (2018;
WS only), Fernald \& Marchman (2012), Fernald, Marchman \& Weisleder (2013; two
batches of participant identifiers, reconciled) and Fernald, Perfors \& Marchman
(2006; WS only); webCDI administrations for BabyView, with item keys mapped to
Wordbank items; and the SEEDLingS CDI administrations. The item set is the 681 WS
items, grouped into four lexical classes (101, 318, 97 and 165 items) for the
difficulty prior described in S2.1. In total 1,520 item-level administrations enter
the Rasch likelihood; a further 95 administrations from Fernald et al.\ (2013),
available only as WS vocabulary totals, enter through the count likelihood of
S2.1.

[AI generated] Processing speed is the looking-while-listening (LWL) reaction
time. For the five Stanford and SEEDLingS datasets hosted in Peekbank (release
2022.1), reaction time on each distractor-initial trial is the latency of the
first shift from the distractor to the target picture (number of 40-Hz samples
before the shift, plus one, times 25 ms), averaged over trials within an
administration; the measurement model uses the log of this per-administration
mean. The SEEDLingS reaction times were derived from EyeLink fixation reports
(Zhu et al.) with the standard Fernald reaction-time window of 300--1800 ms after
target-word onset and an accuracy window of 367--3500 ms; trials with too little
looking in the window were excluded and one mean reaction time was computed per
session. Altogether 1,374 administrations from 350 children enter the model,
between 2 and 7 per child on average by dataset (Adams et al.\ 6.6, SEEDLingS 5.8,
Fernald \& Marchman 3.5, Fernald et al.\ 2006 2.2, Fernald et al.\ 2013 2.1).

[AI generated] Language input is measured in study-specific units and modelled on
the log scale. The three LENA datasets contribute adult words per hour from
day-long recordings at fixed ages: 16 and 18 months (Adams et al., 2018), 18 and
24 months (Fernald et al., 2013) and monthly from 6 to 17 months (SEEDLingS), 808
recordings in all. BabyView contributes, for each of 5,739 head-camera videos
longer than 60 s, the number of adult (female and male speaker) alphabetic tokens
in the automated transcript at a transcript-confidence threshold of 0.3, divided
by the video's duration in hours. Every recording or video is one observation in
the input measurement model, which has a separate measurement standard deviation
per instrument (LENA vs.\ head camera) and a separate mean per study. Input is
available for 219 children (BabyView 22, SEEDLingS 44, Adams et al.\ 66, Fernald
et al.\ 2013 87) and 166 children have both input and processing data.

[AI generated] [check before use] Sample composition. No participant-level
exclusion on birth status or language delay is applied to the six datasets: the
Adams et al.\ (2018) cohort includes preterm as well as full-term toddlers and the
Fernald \& Marchman (2012) cohort includes late talkers. State the composition
here and adjust the Limitations sentence accordingly.

\subsection*{S1.3 Demographic variables}

[AI generated] Sex is Wordbank's participant sex. Maternal education is Wordbank's
caregiver education level converted to years of schooling (none 0, primary 6,
some secondary 9, secondary 12, some college 14, college 16, some graduate 18,
graduate 20) and standardized within language, so that education effects are per
standard deviation of years.

\section*{S2. Model specification and estimation}

\subsection*{S2.1 The accelerating accumulator}

[AI generated] The item likelihood is $P(y_{ijt} = 1) = \mathrm{logit}^{-1}\!\left(
\xi_i + \log H + \kappa_i \log(t / a_0) - \delta_j\right)$, with $H$ a fixed
constant ($\log H = 5.9$) so that $\xi_i$ is on the scale of a log input rate, and
$a_0 = 21$ months in the input/processing fits. The population acceleration is
parametrized as $\kappa = 1 + \delta$ with $\delta \sim N(0, 10)$; child effects
are non-centred. Item difficulties have a prior by lexical class, $\delta_j \sim
N(\mu_{c(j)}, \tau_{c(j)})$ with $\mu_c \sim N(8, 3)$ and $\tau_c \sim
N^+(0, 1)$; collapsing to a single class changes nothing of substance.

[AI generated] Administrations available only as vocabulary totals enter through
a count likelihood: the number of words produced is modelled as a Poisson-binomial
sum over the form's 681 items, approximated by a normal distribution with mean
$\sum_j p_{ij}$ and variance $\sum_j p_{ij}(1 - p_{ij})$ evaluated at the child's
$\theta_i(t)$.

\subsection*{S2.2 Measurement models and structural links}

[AI generated] Input: each recording $v$ of child $i$ in study $s$ gives $r_{iv}
\sim N(\mu_{r,s} + d_i, \sigma_{m(v)})$, with study means $\mu_{r,s} \sim
N(\bar r_s, 1)$ centred on the observed study mean, instrument-specific measurement
standard deviations $\sigma_m \sim N^+(0, 1)$, and child deviations $d_i \sim N(0,
\sigma_r)$ with $\sigma_r \sim N(0.44, 0.1)$. Processing: each LWL administration
$n$ gives $q_{in} \sim N(\tau_s + w_i + \psi \log(t_n / a_0), \sigma_q)$, with
study intercepts $\tau_s \sim N(6.84, 0.2)$ on the log-ms scale, one global
log-age slope $\psi \sim N(-0.35, 0.1)$, child deviations $w_i \sim N(0,
\sigma_w)$ with $\sigma_w \sim N(0.143, 0.04)$, and $\sigma_q \sim N(0.24,
0.05)$.

[AI generated] Structural links: $\xi_i = \mu_\xi + d_i + \beta_w w_i +
\alpha_i$ and $\kappa_i = 1 + \delta + \gamma d_i + \beta'_w w_i + \zeta_i$, with
independent residuals $\alpha_i \sim N(0, \sigma_\alpha)$, $\zeta_i \sim N(0,
\sigma_\zeta)$, $\sigma_\alpha, \sigma_\zeta \sim N^+(0, 1)$; no residual
correlation between $\alpha$ and $\zeta$ is estimated, the two measured channels
carrying the covariance between efficiency and acceleration. The coefficient of
$d_i$ in $\xi_i$ is fixed at 1 because input enters the accumulator as part of the
log rate. Coupling priors are scaled to reference standard deviations of the
predictors, $\gamma \sim N(0, 1/0.44)$ and $\beta_w, \beta'_w \sim N(0,
1/0.143)$, so that each prior implies a unit-logit effect per reference SD.

[AI generated] Reported effects are per standard deviation of the predictor:
input~$\to$~efficiency equals $\sigma_r$ (a consequence of the fixed unit
coefficient, not an estimated coupling), input~$\to$~acceleration equals $\gamma
\sigma_r$, processing~$\to$~efficiency equals $-\beta_w \sigma_w$ and
processing~$\to$~acceleration equals $-\beta'_w \sigma_w$, the sign flipped so that
faster processing is positive. Efficiency effects are in logits; acceleration
effects are in units of $\kappa$ (logits per unit of log age) and convert to
level-equivalent logits at 30 months by the factor $\log(30/21) = 0.36$.

\subsection*{S2.3 Estimation}

[AI generated] The three arms are separate Stan programs on the same data: the
joint model with both channels, an input-only program without the processing
channel, and a processing-only program without the input channel. Couplings were
introduced along a ladder (input coupling only; adding the processing-to-efficiency
coupling; adding the processing-to-acceleration coupling); the reported joint and
processing-only fits are the top rung and the input-only fit has the input coupling
only. Models were fit with CmdStan 2.38.0 through cmdstanr 0.9.0 with 4 chains of
1,000 warm-up and 1,000 sampling iterations, target acceptance 0.95, maximum tree
depth 10, within-chain parallelism, and initial values at the prior means of the
measurement parameters. Point estimates are posterior means and intervals are
2.5/97.5\% posterior quantiles.

[AI generated] Convergence, over the structural and population parameters: joint
model maximum $\hat R$ 1.03 and minimum bulk effective sample size 163; input-only
1.03 and 88; processing-only 1.04 and 112.

[AI generated] Estimates per standard deviation with 95\% intervals, joint model
versus each channel fit alone: input~$\to$~efficiency 0.35 [0.31, 0.40] versus 0.35
[0.31, 0.40]; input~$\to$~acceleration 0.72 [0.02, 1.42] versus 0.76 [$-$0.11,
1.55]; processing~$\to$~efficiency 0.76 [0.46, 1.04] versus 0.82 [0.53, 1.09];
processing~$\to$~acceleration 0.38 [$-$0.22, 0.99] versus 0.50 [$-$0.21, 1.07].
Population parameters agree to two decimals across arms ($\delta$ 10.29 / 10.30 /
10.28; $\sigma_\xi$ 2.02 in all three).

\subsection*{S2.4 Imputed-input model}

[AI generated] The population-level analysis fits the accelerating accumulator
without measurement channels to the pooled English and the Norwegian Wordbank
longitudinal data, with efficiency decomposed as $\xi_i = \mu_r + r_i +
\alpha_i$, where $r_i \sim N(0, \sigma_r)$ is the child's unobserved log input
deviation and $\alpha_i \sim N(0, \sigma_\alpha)$ the remaining efficiency
deviation, and $\kappa_i = 1 + \delta + \zeta_i$. The word-frequency term of the
companion model is removed (item difficulty is absorbed by $\delta_j$) and the
onset parameter is pinned. Because $r_i$ and $\alpha_i$ are not separately
identified, $\sigma_r$ is fixed at 0.35, 0.44 and 0.58 in three fits per language
and the fit estimates $\sigma_\alpha$; the share of efficiency variance
attributable to input is $1 - \pi_\alpha$ with $\pi_\alpha = \sigma_\alpha^2 /
(\sigma_\alpha^2 + \sigma_r^2)$, which equals $\sigma_r^2 / \sigma_\xi^2$ and is
drawn as the analytic curve using $\sigma_\xi$ from the 0.44 fits (English 1.85,
Norwegian 2.49).

[AI generated] The three values of $\sigma_r$ are, respectively, the direct
estimate from the recordings in the input/processing model (0.35 [0.32, 0.39],
study-centred and therefore a lower bound on population variation), a population
re-anchoring of that estimate to all-adult speech (0.44, the dashed line in
Fig.~5) and an upper value from the literature (0.58); the shaded band runs from
the direct estimate's lower 95\% limit (0.32) to 0.58.

[AI generated] These fits used 4 chains of 1,000 warm-up and 2,000 sampling
iterations at target acceptance 0.97, warm-started from shorter runs. Across the
six fits, $\hat R$ for $\sigma_\xi$ and $\pi_\alpha$ ranged from 1.02 to 1.16 and
bulk effective sample sizes from 19 to 152; Norwegian mixes slowly because of a
posterior ridge between the efficiency and acceleration residual scales
(correlation $-0.23$). The largest $\hat R$ in the full summaries (up to 2.0)
belongs to the prior-pinned onset scale and has no bearing on the reported
quantities. The estimated shares sit on the analytic curve to four decimals:
English 3.6\%, 5.7\% and 9.9\% and Norwegian 2.0\%, 3.1\% and 5.4\% at
$\sigma_r$ = 0.35, 0.44 and 0.58; across the band, English 3.0--9.9\% and
Norwegian 1.7--5.4\%. Other population parameters: English $\delta$ 9.89 and
$\sigma_\zeta$ 6.23; Norwegian $\delta$ 11.80 and $\sigma_\zeta$ 7.45.

\subsection*{S2.5 What $\sigma_r$ means in words per hour, and where the anchors come from}

[AI generated] [main-text / Fig.~5 caption candidate] $\sigma_r$ is the standard
deviation of children's log input rate, so it is most easily read as a
multiplicative spread. At the direct estimate $\sigma_r = 0.35$, one standard
deviation is a factor of 1.4 and the 10th--90th percentile range of input spans a
2.5-fold ratio: at the datasets' reference rate of 1,316 adult words per waking
hour, roughly 840 to 2,060 words per hour. At the upper anchor $\sigma_r = 0.58$
the same range is 630 to 2,770 words per hour, a 4.4-fold ratio
(Table~\ref{tab:si-sigma-r-ranges}); the top axis of Fig.~5 gives these ranges
at each anchor.

\begin{table}[htbp]
\centering\footnotesize
\caption{[AI generated] What each value of $\sigma_r$ (the population SD of log
input rate) implies for the spread of input across children, and the input share
of efficiency variance it yields on the analytic curve of Fig.~5. Percentile
ranges are computed at the reference rate, the geometric mean of the four
datasets' geometric-mean rates (1,316 adult words per waking hour); the
percentile ratio does not depend on the reference. Regenerated by
\texttt{paper/tables/sigma\_r\_tables.R}.}
\label{tab:si-sigma-r-ranges}
% Generated by paper/tables/sigma_r_tables.R (reference rate 1,316 words/h = geometric mean of the four datasets' geometric means) -- do not edit by hand.
\begin{tabular}{lp{5.2cm}rrrrr}
\toprule
$\sigma_r$ & Role & $\times/\div$ per SD & 10th--90th pct.\ (words/h at 1,316) & 90th:10th & Share EN & Share NO \\
\midrule
0.32 & lower edge of the band (direct estimate, lower 95\% limit) & 1.38 & 873--1,983 & 2.3 & 3.0\% & 1.6\% \\
0.35 & direct estimate from this paper's recordings (anchor) & 1.42 & 840--2,061 & 2.5 & 3.6\% & 2.0\% \\
0.44 & raw within-sample SD in this paper's datasets; literature median (main anchor) & 1.55 & 749--2,313 & 3.1 & 5.7\% & 3.1\% \\
0.58 & largest within-sample SD in the evidence (anchor) & 1.79 & 626--2,768 & 4.4 & 9.8\% & 5.4\% \\
0.80 & beyond the band (for scale) & 2.23 & 472--3,669 & 7.8 & 18.7\% & 10.3\% \\
\bottomrule
\end{tabular}

\end{table}

[AI generated] The reference rate and the ranges are in the units of the model's
$r_i$: adult tokens per hour of waking observation, the quantity LENA reports as
adult word count per hour and that the BabyView transcripts yield as adult tokens
per hour of video. Rates in child-directed speech only, or from short visit
recordings, are not comparable (a given household can yield a threefold different
rate depending on which speakers and which hours are counted), which is why every
number below is restricted to the all-adult channel.

[AI generated] The three anchors of Fig.~5 come from three views of the same
quantity. (i) The lowest, 0.35 [0.31, 0.39], is the measurement model's own
estimate of $\sigma_r$ from the 6,547 recordings in the joint input/processing
fit: the SD of children's true log input rate within dataset, after removing
recording-to-recording measurement noise, and therefore a lower bound on the
between-child spread in a broader population. (ii) The main anchor, 0.44, is the
raw SD of children's mean log rate: 0.40 pooled within our four datasets (0.36 to
0.45 by dataset) and 0.45 as the median across eight independent Western
all-adult samples in the literature (Table~\ref{tab:si-sigma-r-evidence}), which
still contain measurement noise. (iii) The upper anchor, 0.58, is the largest
within-sample SD in that evidence, from the low-SES Spanish-speaking sample of
Weisleder \& Fernald (2013). The shaded band of Fig.~5 runs from the direct
estimate's lower 95\% limit (0.32) to 0.58.

\begin{table}[htbp]
\centering\footnotesize
\caption{[AI generated] The evidence behind the $\sigma_r$ anchors: all-adult
tokens per waking hour in Western samples. Literature rows are the entries of the
input-estimation validation set (\texttt{studies/input\_estimation/}) for which a
within-sample SD of log rate could be computed from per-child data; ``pooled''
rows combine the sites above them. The paper's own rows give the raw SD of
children's mean log rate per dataset and, last, the measurement model's
noise-corrected estimate. Regenerated by \texttt{paper/tables/sigma\_r\_tables.R}.}
\label{tab:si-sigma-r-evidence}
% Generated by paper/tables/sigma_r_tables.R from studies/input_estimation/validation_set.csv + the io cache -- do not edit by hand.
\setlength{\tabcolsep}{4pt}
\begin{tabular}{@{}lllllr@{}}
\toprule
Source, sample & Language & $N$ & Method & Words/h & SD(log rate) \\
\midrule
\multicolumn{6}{l}{\textit{Literature: all-adult tokens per waking hour, Western samples}} \\
Soderstrom \& Wittebolle (2013), Daycare & Canadian English & 6 & transcript & 934 & 0.47 \\
Soderstrom \& Wittebolle (2013), Home & Canadian English & 6 & transcript & 1,000 & 0.45 \\
Soderstrom \& Wittebolle (2013), combined & Canadian English & 12 & LENA & 966 & 0.44 \\
Sperry et al. (2019), Alabama & US English & 11 & transcript & 2,934 & 0.47 \\
Sperry et al. (2019), pooled & US English & 42 & transcript & 2,200 & 0.43 \\
Sperry et al. (2019), Daly Park & US English & 7 & transcript & 1,750 & 0.32 \\
Sperry et al. (2019), Indiana & US English & 15 & transcript & 2,287 & 0.42 \\
Sperry et al. (2019), Longwood & US English & 6 & transcript & 2,324 & 0.42 \\
Sperry et al. (2019), South Baltimore & US English & 3 & transcript & 1,517 & 0.44 \\
Weisleder \& Fernald (2013) & US Spanish, low SES & 29 & transcript & 938 & 0.58 \\
\multicolumn{6}{l}{\textit{This paper: raw SD of children's mean log rate within dataset (pooled 0.40)}} \\
AM2018 & US English & 66 & LENA & 1,550 & 0.37 \\
BabyView & US English & 22 & head-cam & 1,299 & 0.45 \\
FMW2013 & US English & 87 & LENA & 1,354 & 0.43 \\
SEEDLingS & US English & 44 & LENA & 1,101 & 0.36 \\
\multicolumn{6}{l}{\textit{This paper: measurement model, noise-corrected (95\% posterior interval)}} \\
Joint io/proc fit, four datasets & US English & 219 & LENA + head-cam & -- & 0.35 [0.31, 0.39] \\
\bottomrule
\end{tabular}

\end{table}

[AI generated] The input-estimation validation set collects 54 estimates of
children's language input from the literature and from this paper's recordings,
each converted to adult tokens per hour of waking observation (with a
0.3-second-per-word conversion where a source reports minutes of speech), and
classified by which speech is counted: mother-directed, any-adult child-directed,
all-adult (child-directed and overheard), overheard only, or all input including
overlapping speakers. Sources span Hart \& Risley (1995), Sperry et al.\ (2019),
Weisleder \& Fernald (2013), Soderstrom \& Wittebolle (2013), Gilkerson et al.\
(2017), Bergelson et al.\ (2018, 2019, 2023), Dailey \& Bergelson (2022), Casillas
et al.\ (2020, 2021), Bunce et al.\ (2024), Cristia (2023) and Coffey et al.\
(2024). Across Western samples the SD of log rate within a sample is 0.32--0.47
for all-adult counts and 0.35--0.88 for child-directed counts, and the SD of
sample means across all samples including non-Western communities exceeds 1.0,
which is why $\sigma_r$ values near 1 would require invoking cross-cultural
variation that the Wordbank samples do not contain. [check: cite the validation
set's README for these summary ranges before use.]

\section*{S3. Demographic analyses}

\subsection*{S3.1 Longitudinal datasets}

[AI generated] For each longitudinal dataset (Thal, Marchman, Norwegian,
Japanese) we fit the accelerating accumulator as a generalized linear mixed model
with \texttt{lme4} (version 1.1.37): \texttt{produces $\sim$ log\_age + (1 +
log\_age | child) + (1 | item)}, binomial family, \texttt{nAGQ = 0}, with
\texttt{log\_age} $= \log(t / a_0)$ and $a_0$ the dataset's median administration
age. The child random intercept is the efficiency estimate (ability at $a_0$) and
the child random slope the acceleration estimate; both are taken as conditional
modes (BLUPs). Within each dataset we regressed each BLUP on sex (female minus
male) and on standardized years of maternal education with ordinary least
squares. Sex is available in all four datasets and education in Marchman and
Norwegian. Figure~2A--B pools the two English datasets by inverse-variance
weighting; the longitudinal meta-analytic estimates in Figure~2C are
random-effects (REML) meta-analyses over datasets with \texttt{metafor} (version
4.8.0), $k = 4$ for sex and $k = 2$ for education.

\subsection*{S3.2 Cross-sectional languages}

[AI generated] For each language we fit, on one randomly selected administration
per child, the item-level Rasch GLMM \texttt{produces $\sim$ predictor $\times$
log\_age + (1 | item) + (1 | child)} with \texttt{lme4} (\texttt{nAGQ = 0}), where
\texttt{log\_age} $= \log(t / a_0)$ and $a_0$ is the median age of the language's
sample. The predictor's main effect is its effect on efficiency (ability at
$a_0$) and its interaction with log age its effect on acceleration; sex effects are
reported as female minus male.

[AI generated] Languages were monolingual, typically developing Wordbank samples
(bilingual-origin and clinical datasets excluded) with at least 200 children with
recorded sex; English (British) was excluded (a twin study administered with a
short form) and French (Qu\'ebec) as bilingual, leaving 31 languages. The
education model additionally required at least 1,000 item responses with education
recorded and at least two education levels. Figure~2 displays languages with at
least 300 children (28 for sex, 12 for education); French (France) is omitted
from the education panel because its skewed education distribution gives an
unstable estimate. Cross-language meta-analytic estimates are random-effects
(REML) meta-analyses over the displayed languages with \texttt{metafor}, and the
shaded bands in Figure~2A--B are their 95\% confidence intervals.

[AI generated] [check before use] In an earlier validation on the languages with
both designs, cross-sectional acceleration effects ran 1.5--2.3 times the
corresponding longitudinal effects, so cross-sectional acceleration estimates
should be read as upper bounds.

\section*{S4. Figure conventions}

[AI generated] Figure~1B uses the joint input/processing fit's population
parameters ($\kappa = 11.3$, $\sigma_\xi = 2.02$, $\sigma_\kappa = 4.40$, $a_0
= 21$ months, $\log H = 5.9$) for a reference child 0.5 SD below the mean
efficiency and adds one population SD of efficiency or of acceleration; words
produced are obtained by passing $\theta$ through an illustrative bank of 680
items with logistic difficulties of mean 15 and scale 2.5.

[AI generated] In Figure~2A--B, acceleration confidence intervals are truncated at
$\pm$(largest absolute estimate + 0.7) so that a few small-language estimates do
not compress the axis; the meta-analytic panel (Figure~2C) uses its own scale,
shared by efficiency and acceleration. Figure~3 draws one line per child (BabyView
videos averaged within calendar month; LWL administrations as session means on a
log axis; input in study-specific units) and a loess smoother per dataset (span 1).
Intervals in Figures~4 and~5 are 95\% posterior intervals; the dashed line in
Figure~5 marks $\sigma_r = 0.44$.

\section*{S5. Software and reproducibility}

[AI generated] Analyses used R 4.5.1 with cmdstanr 0.9.0 and CmdStan 2.38.0,
lme4 1.1.37, metafor 4.8.0 and wordbankr 2.0.0; LWL data were drawn from Peekbank
release 2022.1. Code, the raw input and processing data files, the cached model
summaries behind every figure and table, and the scripts that regenerate them are
archived at \url{https://github.com/langcog/standard-model-2}; cross-linguistic
MB-CDI data are downloaded from Wordbank by \texttt{wordbankr} at build time.
